% القسم: المناهج
% الفصل: الاستقراء
% المبحث: المقدمة

\documentclass[../../../include/open-logic-section]{subfiles}

\begin{document}

\olfileid{mth}{ind}{int}

\olsection{تمهيد}

الاستقراء من أمهات طرق البرهان، وتختلف صوره فيكاد لا يخلو منه باب من
أبواب المنطق وعلم الحاسوب النظري والرياضيات. وستدعونا الحاجة إليه في
إثبات نتائج منطقية كثيرة.

ويُجعل الاستقراء في مقابلة الاستنباط، ويقال إنه انتقال من الجزئي إلى
الكلي. فمن رأى زمردات كثيرة كلها خضر، ولم ير ما يسمى زمردًا وهو غير
أخضر، ربما حكم بأن الزمرد كله أخضر. وهذا استقراء لأنه جمع أحكامًا على
جزئيات كثيرة---هذه خضراء وتلك خضراء وما أشبههما---ثم انتهى إلى الحكم
العام بأن كل زمرد أخضر. وأما الاستقراء \emph{الرياضي} فينتهي هو أيضًا
إلى قضية عامة، غير أن بينه وبين هذا «الاستقراء البسيط» فرقًا عظيمًا.

وحاصل البرهان الاستقرائي الرياضي، على سبيل التقريب، أن جميع الأشياء
الرياضية من صنف معلوم تشترك في خاصية. وأيسر مواضعه أن تكون تلك
الأشياء أعدادًا طبيعية؛ فلكي نثبت الخاصية لكل عدد طبيعي نبيّن أمرين:
\begin{enumerate}
    \item أن العدد $0$ متصف بالخاصية، و
    \item أن اتصاف العدد~$k$ بها يستلزم اتصاف العدد~$k+1$ بها أيضاً.
\end{enumerate}
وقد يمتد الاستقراء على الأعداد الطبيعية إلى قضايا عامة في أشياء
رياضية يمكن إسناد عدد إلى كل واحد منها. فالمجموعة المتناهية، مثلًا،
لها عدد متناه~$n$ من العناصر؛ فإذا برهنا بالاستقراء، لكل عدد~$n$، أن
«جميع المجموعات المتناهية التي حجمها~$n$ هي \dots»، فقد حصل لنا حكم
يشمل جميع المجموعات المتناهية.

ثم يعمّم الاستقراء إلى الأشياء الرياضية \emph{المعرّفة استقرائيًّا}.
فعبارات اللغة الصورية، ومنها عبارات منطق الرتبة الأولى، تُعرّف بهذا
الوجه. والطريق إلى إثبات حكم يعم جميع تلك العبارات هو
\emph{الاستقراء البنيوي}، وهو من أنفع الأدوات وأكثرها دورانًا في
المنطق.

\end{document}
